\QHeading{第1题　化简下面的小数。}
\par\noindent


\begin{oralgrid}
\OralItem{1}{1.60}{1.6}
\OralItem{2}{0.80}{0.8}
\OralItem{3}{10.060}{10.06}
\OralItem{4}{80.00}{80}
\OralItem{5}{5.050}{5.05}
\OralItem{6}{0.900}{0.9}
\OralItem{7}{30.70}{30.7}
\OralItem{8}{0.0320}{0.032}
\OralItem{9}{70.040}{70.04}
\OralItem{10}{0.090}{0.09}
\OralItem{11}{13.090}{13.09}
\OralItem{12}{0.7040}{0.704}\end{oralgrid}
\par

\par\noindent

\QHeading{第2题　不改变数的大小，把下列各数改写成两位小数。}

\PlainItem{$10=$（{\color{answerred}$10.00$}）}
\par

\PlainItem{$3.200=$（{\color{answerred}$3.20$}）}
\PlainItem{$21.8=$（{\color{answerred}$21.80$}）}
\PlainItem{$0.7=$（{\color{answerred}$0.70$}）}
\PlainItem{$1.7=$（{\color{answerred}$1.70$}）}
\PlainItem{$20.700=$（{\color{answerred}$20.70$}）}
\PlainItem{$10.6=$（{\color{answerred}$10.60$}）}
\PlainItem{$6.9=$（{\color{answerred}$6.90$}）}
\PlainItem{$1.050=$（{\color{answerred}$1.05$}）}
\par\noindent

\QHeading{第3题　把下面物品的价格改写成用“元”作单位的两位小数。}

\par\noindent\makebox[\linewidth][c]{
\begin{tikzpicture}[x=1bp,y=-1bp,color=black,draw=black,text=black]
\begin{scope}\clip (67.0,392.0) rectangle (543.0,522.0);
\hdhead{49}{小数的意义与性质——信息窗2}
\hdsection{1}{154}{化简下面的小数。}
\q{109}{201}{185}{13}{$1.60=$}
\q{109}{246}{185}{13}{$5.050=$}
\q{109}{291}{185}{13}{$70.040=$}
\q{304}{202}{185}{13}{$0.80=$}
\q{304}{246}{185}{13}{$0.900=$}
\q{304}{291}{185}{13}{$0.090=$}
\q{500}{203}{185}{13}{$10.060=$}
\q{500}{247}{185}{13}{$30.70=$}
\q{500}{292}{185}{13}{$13.090=$}
\q{691}{202}{185}{13}{$80.00=$}
\q{691}{247}{185}{13}{$0.0320=$}
\q{691}{292}{185}{13}{$0.7040=$}
\hdsection{2}{352}{不改变数的大小，把下列各数改写成两位小数。}
\q{109}{398}{190}{13}{$10=$}
\q{155}{398}{28}{13}{（}
\q{229}{398}{25}{13}{）}
\q{372}{398}{190}{13}{$3.200=$}
\q{442}{398}{28}{13}{（}
\q{523}{398}{25}{13}{）}
\q{627}{398}{190}{13}{$21.8=$}
\q{699}{398}{28}{13}{（}
\q{768}{398}{25}{13}{）}
\q{109}{444}{190}{13}{$0.7=$}
\q{161}{444}{28}{13}{（}
\q{239}{444}{25}{13}{）}
\q{372}{444}{190}{13}{$1.7=$}
\q{420}{444}{28}{13}{（}
\q{501}{444}{25}{13}{）}
\q{627}{444}{190}{13}{$20.700=$}
\q{720}{444}{28}{13}{（}
\q{795}{444}{25}{13}{）}
\q{109}{490}{190}{13}{$10.6=$}
\q{172}{490}{28}{13}{（}
\q{252}{490}{25}{13}{）}
\q{372}{490}{190}{13}{$6.9=$}
\q{420}{490}{28}{13}{（}
\q{501}{490}{25}{13}{）}
\q{627}{490}{190}{13}{$1.050=$}
\q{711}{490}{28}{13}{（}
\q{787}{490}{25}{13}{）}
\hdsection{3}{547}{把下面物品的价格改写成用“元”作单位的两位小数。}
\begin{scope}[scale=.6666667,draw=black,text=black]\node[inner sep=0pt,font=\normalsize,text=black] at (119.000,608.000) {①};\end{scope}
\begin{scope}[scale=.6666667,draw=black,text=black]\node[inner sep=0pt,font=\normalsize,text=black] at (311.000,608.000) {②};\end{scope}
\begin{scope}[scale=.6666667,draw=black,text=black]\node[inner sep=0pt,font=\normalsize,text=black] at (500.000,608.000) {③};\end{scope}
\begin{scope}[scale=.6666667,draw=black,text=black]\node[inner sep=0pt,font=\normalsize,text=black] at (685.000,608.000) {④};\end{scope}
\begin{scope}[scale=.6666667,draw=black,text=black]\draw[draw=black,line width=.7pt,fill=bookpink!45] (154.000,671.000)--(166.000,610.000)--(219.000,611.000)--(214.000,673.000)--cycle;\end{scope}
\begin{scope}[scale=.6666667,draw=black,text=black]\draw[draw=black!45,line width=.3pt] (160.000,667.000)--(171.000,614.000)--(213.000,615.000)--(209.000,669.000)--cycle;\end{scope}
\begin{scope}[scale=.6666667,draw=black,text=black]\draw[line width=.4pt] (173,619) rectangle (206,630);\draw[line width=.5pt] (174,655)--(198,655) (173,660)--(197,660);\end{scope}
\begin{scope}[scale=.6666667,draw=black,text=black]\draw[draw=black,line width=.6pt,fill=yellow!18] (347.000,625.000)--(350.000,620.000)--(421.000,643.000)--(433.000,651.000)--(419.000,647.000)--cycle;\end{scope}
\begin{scope}[scale=.6666667,draw=black,text=black]\draw[draw=black,line width=.4pt] (354.000,624.000)--(420.000,645.000);\end{scope}
\begin{scope}[scale=.6666667,draw=black,text=black]\draw[draw=black,line width=.3pt,fill=black] (426.000,646.000)--(433.000,651.000)--(425.000,649.000)--cycle;\end{scope}
\begin{scope}[scale=.6666667,draw=black,text=black]\draw[draw=black,line width=.55pt,fill=bookpink] (540.000,664.000)--(550.000,625.000)--(584.000,626.000)--(603.000,667.000)--cycle;\end{scope}
\begin{scope}[scale=.6666667,draw=black,text=black]\draw[draw=black,line width=.35pt] (541.000,666.000)--(603.000,669.000);\end{scope}
\begin{scope}[scale=.6666667,draw=black,text=black]\draw[draw=black,line width=.55pt,fill=bookpink] (540.000,659.000)--(550.000,620.000)--(584.000,621.000)--(603.000,662.000)--cycle;\end{scope}
\begin{scope}[scale=.6666667,draw=black,text=black]\draw[draw=black,line width=.35pt] (541.000,661.000)--(603.000,664.000);\end{scope}
\begin{scope}[scale=.6666667,draw=black,text=black]\draw[draw=black,line width=.55pt,fill=bookpink] (540.000,655.000)--(550.000,616.000)--(584.000,617.000)--(603.000,658.000)--cycle;\end{scope}
\begin{scope}[scale=.6666667,draw=black,text=black]\draw[draw=black,line width=.35pt] (541.000,657.000)--(603.000,660.000);\end{scope}
\begin{scope}[scale=.6666667,draw=black,text=black]\draw[line width=.6pt] (547.0,620) ellipse[x radius=5,y radius=3];\end{scope}
\begin{scope}[scale=.6666667,draw=black,text=black]\draw[line width=.6pt] (548.8,628) ellipse[x radius=5,y radius=3];\end{scope}
\begin{scope}[scale=.6666667,draw=black,text=black]\draw[line width=.6pt] (550.6,636) ellipse[x radius=5,y radius=3];\end{scope}
\begin{scope}[scale=.6666667,draw=black,text=black]\draw[line width=.6pt] (552.4,644) ellipse[x radius=5,y radius=3];\end{scope}
\begin{scope}[scale=.6666667,draw=black,text=black]\draw[line width=.6pt] (554.2,652) ellipse[x radius=5,y radius=3];\end{scope}
\begin{scope}[scale=.6666667,draw=black,text=black]\draw[draw=black,line width=.7pt] (561.000,630.000)--(581.000,631.000);\end{scope}
\begin{scope}[scale=.6666667,draw=black,text=black]\draw[draw=black,line width=.7pt] (560.000,638.000)--(586.000,639.000);\end{scope}
\begin{scope}[scale=.6666667,draw=black,text=black]\path[draw=black,line width=.65pt,fill=bookpink,rounded corners=4] (722,603) rectangle (774,680);\path[draw=black,line width=.5pt,fill=white,rounded corners=2] (726,609) rectangle (770,624);\end{scope}
\begin{scope}[scale=.6666667,draw=black,text=black]\node[inner sep=0pt,font=\normalsize,text=black] at (761.000,616.000) {0};\end{scope}
\begin{scope}[scale=.6666667,draw=black,text=black]\path[draw=black,line width=.4pt,fill=white,rounded corners=1] (726,630) rectangle (734,636);\end{scope}
\begin{scope}[scale=.6666667,draw=black,text=black]\path[draw=black,line width=.4pt,fill=white,rounded corners=1] (726,641) rectangle (734,647);\end{scope}
\begin{scope}[scale=.6666667,draw=black,text=black]\path[draw=black,line width=.4pt,fill=white,rounded corners=1] (726,652) rectangle (734,658);\end{scope}
\begin{scope}[scale=.6666667,draw=black,text=black]\path[draw=black,line width=.4pt,fill=white,rounded corners=1] (726,663) rectangle (734,669);\end{scope}
\begin{scope}[scale=.6666667,draw=black,text=black]\path[draw=black,line width=.4pt,fill=white,rounded corners=1] (737,630) rectangle (745,636);\end{scope}
\begin{scope}[scale=.6666667,draw=black,text=black]\path[draw=black,line width=.4pt,fill=white,rounded corners=1] (737,641) rectangle (745,647);\end{scope}
\begin{scope}[scale=.6666667,draw=black,text=black]\path[draw=black,line width=.4pt,fill=white,rounded corners=1] (737,652) rectangle (745,658);\end{scope}
\begin{scope}[scale=.6666667,draw=black,text=black]\path[draw=black,line width=.4pt,fill=white,rounded corners=1] (737,663) rectangle (745,669);\end{scope}
\begin{scope}[scale=.6666667,draw=black,text=black]\path[draw=black,line width=.4pt,fill=white,rounded corners=1] (748,630) rectangle (756,636);\end{scope}
\begin{scope}[scale=.6666667,draw=black,text=black]\path[draw=black,line width=.4pt,fill=white,rounded corners=1] (748,641) rectangle (756,647);\end{scope}
\begin{scope}[scale=.6666667,draw=black,text=black]\path[draw=black,line width=.4pt,fill=white,rounded corners=1] (748,652) rectangle (756,658);\end{scope}
\begin{scope}[scale=.6666667,draw=black,text=black]\path[draw=black,line width=.4pt,fill=white,rounded corners=1] (748,663) rectangle (756,669);\end{scope}
\begin{scope}[scale=.6666667,draw=black,text=black]\path[draw=black,line width=.4pt,fill=white,rounded corners=1] (759,630) rectangle (767,636);\end{scope}
\begin{scope}[scale=.6666667,draw=black,text=black]\path[draw=black,line width=.4pt,fill=white,rounded corners=1] (759,641) rectangle (767,647);\end{scope}
\begin{scope}[scale=.6666667,draw=black,text=black]\path[draw=black,line width=.4pt,fill=white,rounded corners=1] (759,652) rectangle (767,658);\end{scope}
\begin{scope}[scale=.6666667,draw=black,text=black]\path[draw=black,line width=.4pt,fill=white,rounded corners=1] (759,663) rectangle (767,669);\end{scope}
\q{153}{708}{155}{13}{2元5角}
\q{371}{708}{155}{13}{7角}
\q{533}{708}{155}{13}{6元8分}
\q{704}{708}{155}{13}{35元5角}
\q{143}{756}{28}{13}{（}
\q{214}{756}{25}{13}{）}
\q{234}{756}{45}{13}{元}
\q{337}{756}{28}{13}{（}
\q{405}{756}{25}{13}{）}
\q{425}{756}{45}{13}{元}
\q{524}{756}{28}{13}{（}
\q{590}{756}{25}{13}{）}
\q{610}{756}{45}{13}{元}
\q{699}{756}{28}{13}{（}
\q{768}{756}{25}{13}{）}
\q{788}{756}{45}{13}{元}
\hdsection{4}{816}{解决问题。}
\q{109}{866}{734}{12.5}{6名同学参加50米短跑比赛，成绩如下：}
\begin{scope}[scale=.6666667,draw=black,text=black]\path[pattern=north east lines,pattern color=black] (190,909) rectangle (317,987);\end{scope}
\begin{scope}[scale=.6666667,draw=black,text=black]\draw[draw=black,line width=.5pt] (190.000,909.000)--(190.000,987.000);\end{scope}
\begin{scope}[scale=.6666667,draw=black,text=black]\draw[draw=black,line width=.5pt] (317.000,909.000)--(317.000,987.000);\end{scope}
\begin{scope}[scale=.6666667,draw=black,text=black]\draw[draw=black,line width=.5pt] (395.000,909.000)--(395.000,987.000);\end{scope}
\begin{scope}[scale=.6666667,draw=black,text=black]\draw[draw=black,line width=.5pt] (472.000,909.000)--(472.000,987.000);\end{scope}
\begin{scope}[scale=.6666667,draw=black,text=black]\draw[draw=black,line width=.5pt] (547.000,909.000)--(547.000,987.000);\end{scope}
\begin{scope}[scale=.6666667,draw=black,text=black]\draw[draw=black,line width=.5pt] (623.000,909.000)--(623.000,987.000);\end{scope}
\begin{scope}[scale=.6666667,draw=black,text=black]\draw[draw=black,line width=.5pt] (699.000,909.000)--(699.000,987.000);\end{scope}
\begin{scope}[scale=.6666667,draw=black,text=black]\draw[draw=black,line width=.5pt] (774.000,909.000)--(774.000,987.000);\end{scope}
\begin{scope}[scale=.6666667,draw=black,text=black]\draw[draw=black,line width=.5pt] (190.000,909.000)--(774.000,909.000);\end{scope}
\begin{scope}[scale=.6666667,draw=black,text=black]\draw[draw=black,line width=.5pt] (190.000,948.000)--(774.000,948.000);\end{scope}
\begin{scope}[scale=.6666667,draw=black,text=black]\draw[draw=black,line width=.5pt] (190.000,987.000)--(774.000,987.000);\end{scope}
\begin{scope}[scale=.6666667,draw=black,text=black]\node[inner sep=0pt,font=\normalsize,text=black] at (253.500,928.500) {姓名};\end{scope}
\begin{scope}[scale=.6666667,draw=black,text=black]\node[inner sep=0pt,font=\normalsize,text=black] at (356.000,928.500) {小强};\end{scope}
\begin{scope}[scale=.6666667,draw=black,text=black]\node[inner sep=0pt,font=\normalsize,text=black] at (433.500,928.500) {小东};\end{scope}
\begin{scope}[scale=.6666667,draw=black,text=black]\node[inner sep=0pt,font=\normalsize,text=black] at (509.500,928.500) {小刚};\end{scope}
\begin{scope}[scale=.6666667,draw=black,text=black]\node[inner sep=0pt,font=\normalsize,text=black] at (585.000,928.500) {小庆};\end{scope}
\begin{scope}[scale=.6666667,draw=black,text=black]\node[inner sep=0pt,font=\normalsize,text=black] at (661.000,928.500) {小杰};\end{scope}
\begin{scope}[scale=.6666667,draw=black,text=black]\node[inner sep=0pt,font=\normalsize,text=black] at (736.500,928.500) {小涛};\end{scope}
\begin{scope}[scale=.6666667,draw=black,text=black]\node[inner sep=0pt,font=\normalsize,text=black] at (253.500,967.500) {成绩/秒};\end{scope}
\begin{scope}[scale=.6666667,draw=black,text=black]\node[inner sep=0pt,font=\normalsize,text=black] at (356.000,967.500) {8.9};\end{scope}
\begin{scope}[scale=.6666667,draw=black,text=black]\node[inner sep=0pt,font=\normalsize,text=black] at (433.500,967.500) {8.1};\end{scope}
\begin{scope}[scale=.6666667,draw=black,text=black]\node[inner sep=0pt,font=\normalsize,text=black] at (509.500,967.500) {9.6};\end{scope}
\begin{scope}[scale=.6666667,draw=black,text=black]\node[inner sep=0pt,font=\normalsize,text=black] at (585.000,967.500) {8.3};\end{scope}
\begin{scope}[scale=.6666667,draw=black,text=black]\node[inner sep=0pt,font=\normalsize,text=black] at (661.000,967.500) {9.0};\end{scope}
\begin{scope}[scale=.6666667,draw=black,text=black]\node[inner sep=0pt,font=\normalsize,text=black] at (736.500,967.500) {8.7};\end{scope}
\q{109}{1012}{700}{12.5}{请把前三名的名字写在领奖台上。}
\begin{scope}[scale=.6666667,draw=black,text=black]\draw[draw=black,line width=.6pt,fill=bookpink] (397.000,1140.000)--(420.000,1127.000)--(441.000,1127.000)--(441.000,1140.000)--cycle;\end{scope}
\begin{scope}[scale=.6666667,draw=black,text=black]\draw[draw=black,line width=.6pt,fill=bookpink] (441.000,1116.000)--(465.000,1102.000)--(510.000,1102.000)--(486.000,1116.000)--cycle;\end{scope}
\begin{scope}[scale=.6666667,draw=black,text=black]\draw[draw=black,line width=.6pt,fill=bookpink] (486.000,1148.000)--(510.000,1134.000)--(554.000,1134.000)--(529.000,1148.000)--cycle;\end{scope}
\begin{scope}[scale=.6666667,draw=black,text=black]\draw[draw=black,line width=.6pt,fill=bookpink!75!black] (486.000,1116.000)--(510.000,1102.000)--(510.000,1134.000)--(486.000,1148.000)--cycle;\end{scope}
\begin{scope}[scale=.6666667,draw=black,text=black]\draw[draw=black,line width=.6pt,fill=bookpink!75!black] (529.000,1148.000)--(554.000,1134.000)--(554.000,1157.000)--(529.000,1171.000)--cycle;\end{scope}
\begin{scope}[scale=.6666667,draw=black,text=black]\draw[draw=black,line width=.6pt,fill=bookpink] (397.000,1140.000)--(441.000,1140.000)--(441.000,1116.000)--(486.000,1116.000)--(486.000,1148.000)--(529.000,1148.000)--(529.000,1171.000)--(397.000,1171.000)--cycle;\end{scope}
\begin{scope}[scale=.6666667,draw=black,text=black]\node[inner sep=0pt,font=\normalsize,text=black] at (464.000,1133.000) {1};\end{scope}
\begin{scope}[scale=.6666667,draw=black,text=black]\node[inner sep=0pt,font=\normalsize,text=black] at (417.000,1156.000) {2};\end{scope}
\begin{scope}[scale=.6666667,draw=black,text=black]\node[inner sep=0pt,font=\normalsize,text=black] at (506.000,1159.000) {3};\end{scope}

\begin{scope}[color=answerred,draw=answerred,text=answerred]\ans{114.00}{134.00}{63.33}{9.5}{$\vphantom{0}1.6$}
\ans{121.33}{164.00}{63.33}{9.5}{$\vphantom{0}5.05$}
\ans{130.67}{194.00}{63.33}{9.5}{$\vphantom{0}70.04$}
\ans{247.33}{134.67}{63.33}{9.5}{$\vphantom{0}0.8$}
\ans{254.67}{164.00}{63.33}{9.5}{$\vphantom{0}0.9$}
\ans{254.67}{194.00}{63.33}{9.5}{$\vphantom{0}0.09$}
\ans{391.33}{135.33}{63.33}{9.5}{$\vphantom{0}10.06$}
\ans{382.00}{164.67}{63.33}{9.5}{$\vphantom{0}30.7$}
\ans{391.33}{194.67}{63.33}{9.5}{$\vphantom{0}13.09$}
\ans{512.67}{134.67}{63.33}{9.5}{$\vphantom{0}80$}
\ans{523.33}{164.67}{63.33}{9.5}{$\vphantom{0}0.032$}
\ans{523.33}{194.67}{63.33}{9.5}{$\vphantom{0}0.704$}
\ans{118.00}{264.67}{66.67}{10}{$\vphantom{0}10.00$}
\ans{309.33}{264.67}{66.67}{10}{$\vphantom{0}3.20$}
\ans{480.67}{264.67}{66.67}{10}{$\vphantom{0}21.80$}
\ans{122.00}{295.33}{66.67}{10}{$\vphantom{0}0.70$}
\ans{294.67}{295.33}{66.67}{10}{$\vphantom{0}1.70$}
\ans{494.67}{295.33}{66.67}{10}{$\vphantom{0}20.70$}
\ans{129.33}{326.00}{66.67}{10}{$\vphantom{0}10.60$}
\ans{294.67}{326.00}{66.67}{10}{$\vphantom{0}6.90$}
\ans{488.67}{326.00}{66.67}{10}{$\vphantom{0}1.05$}
\ans{109.33}{503.33}{58.67}{11}{$\vphantom{0}2.50$}
\ans{238.00}{503.33}{58.67}{11}{$\vphantom{0}0.70$}
\ans{362.67}{503.33}{58.67}{11}{$\vphantom{0}6.08$}
\ans{479.33}{503.33}{58.67}{11}{$\vphantom{0}35.50$}
\ans{72.67}{696.67}{467.33}{10}{用时越少，名次越靠前：$8.1<8.3<8.7<8.9<9.0<9.6$。}
\ans{294.00}{719.33}{58.67}{10}{小东}
\ans{244.00}{748.00}{48.00}{9.5}{小庆}
\ans{363.33}{753.33}{48.00}{9.5}{小涛}
\draw[answerred,line width=.7pt] (270.000,757.333) -- (280.000,763.333);
\draw[answerred,line width=.7pt] (365.333,764.667) -- (346.667,766.667);

\ans{72.67}{792}{470}{10}{答：第一名是小东，第二名是小庆，第三名是小涛。}
\end{scope}
\end{scope}\pgfresetboundingbox\path[use as bounding box] (67.0,392.0) rectangle (543.0,522.0);
\end{tikzpicture}}\par



\par






\par\noindent

\Answer{
\begingroup\color{answerred}
\QHeading{计算过程：3. 把价格改写成以元为单位的两位小数}

\noindent
1. \Formula{\begin{aligned}[t]&\phantom{={}}2\text{元}5\text{角}\\&=(2+5\div10)\text{元}\\&=(2+0.5)\text{元}\\&=2.50\text{元}\end{aligned}} \par\medskip
 2. \Formula{\begin{aligned}[t]&\phantom{={}}7\text{角}\\&=(7\div10)\text{元}\\&=0.7\text{元}\\&=0.70\text{元}\end{aligned}}。\par\medskip

3. \Formula{\begin{aligned}[t]&\phantom{={}}6\text{元}8\text{分}\\&=(6+8\div100)\text{元}\\&=(6+0.08)\text{元}\\&=6.08\text{元}\end{aligned}} \par\medskip
 4. \Formula{\begin{aligned}[t]&\phantom{={}}35\text{元}5\text{角}\\&=(35+5\div10)\text{元}\\&=(35+0.5)\text{元}\\&=35.50\text{元}\end{aligned}}



\endgroup

}

\QHeading{第4题　解决问题。}

\Problem{6名同学参加50米短跑比赛，成绩如下：}
\par

\DataTable{2}{姓名 & 成绩/秒\\\hline
小强 & 8.9\\\hline
小东 & 8.1\\\hline
小刚 & 9.6\\\hline
小庆 & 8.3\\\hline
小杰 & 9.0\\\hline
小涛 & 8.7\\\hline}
\Problem{请把前三名的名字写在领奖台上。}
\Answer{用时越少，名次越靠前：$8.1<8.3<8.7<8.9<9.0<9.6$。}



\par\noindent\makebox[\linewidth][c]{
\begin{tikzpicture}[x=1bp,y=-1bp,color=black,draw=black,text=black]
\begin{scope}\clip (230.0,715.0) rectangle (420.0,789.0);
\hdhead{49}{小数的意义与性质——信息窗2}
\hdsection{1}{154}{化简下面的小数。}
\q{109}{201}{185}{13}{$1.60=$}
\q{109}{246}{185}{13}{$5.050=$}
\q{109}{291}{185}{13}{$70.040=$}
\q{304}{202}{185}{13}{$0.80=$}
\q{304}{246}{185}{13}{$0.900=$}
\q{304}{291}{185}{13}{$0.090=$}
\q{500}{203}{185}{13}{$10.060=$}
\q{500}{247}{185}{13}{$30.70=$}
\q{500}{292}{185}{13}{$13.090=$}
\q{691}{202}{185}{13}{$80.00=$}
\q{691}{247}{185}{13}{$0.0320=$}
\q{691}{292}{185}{13}{$0.7040=$}
\hdsection{2}{352}{不改变数的大小，把下列各数改写成两位小数。}
\q{109}{398}{190}{13}{$10=$}
\q{155}{398}{28}{13}{（}
\q{229}{398}{25}{13}{）}
\q{372}{398}{190}{13}{$3.200=$}
\q{442}{398}{28}{13}{（}
\q{523}{398}{25}{13}{）}
\q{627}{398}{190}{13}{$21.8=$}
\q{699}{398}{28}{13}{（}
\q{768}{398}{25}{13}{）}
\q{109}{444}{190}{13}{$0.7=$}
\q{161}{444}{28}{13}{（}
\q{239}{444}{25}{13}{）}
\q{372}{444}{190}{13}{$1.7=$}
\q{420}{444}{28}{13}{（}
\q{501}{444}{25}{13}{）}
\q{627}{444}{190}{13}{$20.700=$}
\q{720}{444}{28}{13}{（}
\q{795}{444}{25}{13}{）}
\q{109}{490}{190}{13}{$10.6=$}
\q{172}{490}{28}{13}{（}
\q{252}{490}{25}{13}{）}
\q{372}{490}{190}{13}{$6.9=$}
\q{420}{490}{28}{13}{（}
\q{501}{490}{25}{13}{）}
\q{627}{490}{190}{13}{$1.050=$}
\q{711}{490}{28}{13}{（}
\q{787}{490}{25}{13}{）}
\hdsection{3}{547}{把下面物品的价格改写成用“元”作单位的两位小数。}
\begin{scope}[scale=.6666667,draw=black,text=black]\node[inner sep=0pt,font=\normalsize,text=black] at (119.000,608.000) {①};\end{scope}
\begin{scope}[scale=.6666667,draw=black,text=black]\node[inner sep=0pt,font=\normalsize,text=black] at (311.000,608.000) {②};\end{scope}
\begin{scope}[scale=.6666667,draw=black,text=black]\node[inner sep=0pt,font=\normalsize,text=black] at (500.000,608.000) {③};\end{scope}
\begin{scope}[scale=.6666667,draw=black,text=black]\node[inner sep=0pt,font=\normalsize,text=black] at (685.000,608.000) {④};\end{scope}
\begin{scope}[scale=.6666667,draw=black,text=black]\draw[draw=black,line width=.7pt,fill=bookpink!45] (154.000,671.000)--(166.000,610.000)--(219.000,611.000)--(214.000,673.000)--cycle;\end{scope}
\begin{scope}[scale=.6666667,draw=black,text=black]\draw[draw=black!45,line width=.3pt] (160.000,667.000)--(171.000,614.000)--(213.000,615.000)--(209.000,669.000)--cycle;\end{scope}
\begin{scope}[scale=.6666667,draw=black,text=black]\draw[line width=.4pt] (173,619) rectangle (206,630);\draw[line width=.5pt] (174,655)--(198,655) (173,660)--(197,660);\end{scope}
\begin{scope}[scale=.6666667,draw=black,text=black]\draw[draw=black,line width=.6pt,fill=yellow!18] (347.000,625.000)--(350.000,620.000)--(421.000,643.000)--(433.000,651.000)--(419.000,647.000)--cycle;\end{scope}
\begin{scope}[scale=.6666667,draw=black,text=black]\draw[draw=black,line width=.4pt] (354.000,624.000)--(420.000,645.000);\end{scope}
\begin{scope}[scale=.6666667,draw=black,text=black]\draw[draw=black,line width=.3pt,fill=black] (426.000,646.000)--(433.000,651.000)--(425.000,649.000)--cycle;\end{scope}
\begin{scope}[scale=.6666667,draw=black,text=black]\draw[draw=black,line width=.55pt,fill=bookpink] (540.000,664.000)--(550.000,625.000)--(584.000,626.000)--(603.000,667.000)--cycle;\end{scope}
\begin{scope}[scale=.6666667,draw=black,text=black]\draw[draw=black,line width=.35pt] (541.000,666.000)--(603.000,669.000);\end{scope}
\begin{scope}[scale=.6666667,draw=black,text=black]\draw[draw=black,line width=.55pt,fill=bookpink] (540.000,659.000)--(550.000,620.000)--(584.000,621.000)--(603.000,662.000)--cycle;\end{scope}
\begin{scope}[scale=.6666667,draw=black,text=black]\draw[draw=black,line width=.35pt] (541.000,661.000)--(603.000,664.000);\end{scope}
\begin{scope}[scale=.6666667,draw=black,text=black]\draw[draw=black,line width=.55pt,fill=bookpink] (540.000,655.000)--(550.000,616.000)--(584.000,617.000)--(603.000,658.000)--cycle;\end{scope}
\begin{scope}[scale=.6666667,draw=black,text=black]\draw[draw=black,line width=.35pt] (541.000,657.000)--(603.000,660.000);\end{scope}
\begin{scope}[scale=.6666667,draw=black,text=black]\draw[line width=.6pt] (547.0,620) ellipse[x radius=5,y radius=3];\end{scope}
\begin{scope}[scale=.6666667,draw=black,text=black]\draw[line width=.6pt] (548.8,628) ellipse[x radius=5,y radius=3];\end{scope}
\begin{scope}[scale=.6666667,draw=black,text=black]\draw[line width=.6pt] (550.6,636) ellipse[x radius=5,y radius=3];\end{scope}
\begin{scope}[scale=.6666667,draw=black,text=black]\draw[line width=.6pt] (552.4,644) ellipse[x radius=5,y radius=3];\end{scope}
\begin{scope}[scale=.6666667,draw=black,text=black]\draw[line width=.6pt] (554.2,652) ellipse[x radius=5,y radius=3];\end{scope}
\begin{scope}[scale=.6666667,draw=black,text=black]\draw[draw=black,line width=.7pt] (561.000,630.000)--(581.000,631.000);\end{scope}
\begin{scope}[scale=.6666667,draw=black,text=black]\draw[draw=black,line width=.7pt] (560.000,638.000)--(586.000,639.000);\end{scope}
\begin{scope}[scale=.6666667,draw=black,text=black]\path[draw=black,line width=.65pt,fill=bookpink,rounded corners=4] (722,603) rectangle (774,680);\path[draw=black,line width=.5pt,fill=white,rounded corners=2] (726,609) rectangle (770,624);\end{scope}
\begin{scope}[scale=.6666667,draw=black,text=black]\node[inner sep=0pt,font=\normalsize,text=black] at (761.000,616.000) {0};\end{scope}
\begin{scope}[scale=.6666667,draw=black,text=black]\path[draw=black,line width=.4pt,fill=white,rounded corners=1] (726,630) rectangle (734,636);\end{scope}
\begin{scope}[scale=.6666667,draw=black,text=black]\path[draw=black,line width=.4pt,fill=white,rounded corners=1] (726,641) rectangle (734,647);\end{scope}
\begin{scope}[scale=.6666667,draw=black,text=black]\path[draw=black,line width=.4pt,fill=white,rounded corners=1] (726,652) rectangle (734,658);\end{scope}
\begin{scope}[scale=.6666667,draw=black,text=black]\path[draw=black,line width=.4pt,fill=white,rounded corners=1] (726,663) rectangle (734,669);\end{scope}
\begin{scope}[scale=.6666667,draw=black,text=black]\path[draw=black,line width=.4pt,fill=white,rounded corners=1] (737,630) rectangle (745,636);\end{scope}
\begin{scope}[scale=.6666667,draw=black,text=black]\path[draw=black,line width=.4pt,fill=white,rounded corners=1] (737,641) rectangle (745,647);\end{scope}
\begin{scope}[scale=.6666667,draw=black,text=black]\path[draw=black,line width=.4pt,fill=white,rounded corners=1] (737,652) rectangle (745,658);\end{scope}
\begin{scope}[scale=.6666667,draw=black,text=black]\path[draw=black,line width=.4pt,fill=white,rounded corners=1] (737,663) rectangle (745,669);\end{scope}
\begin{scope}[scale=.6666667,draw=black,text=black]\path[draw=black,line width=.4pt,fill=white,rounded corners=1] (748,630) rectangle (756,636);\end{scope}
\begin{scope}[scale=.6666667,draw=black,text=black]\path[draw=black,line width=.4pt,fill=white,rounded corners=1] (748,641) rectangle (756,647);\end{scope}
\begin{scope}[scale=.6666667,draw=black,text=black]\path[draw=black,line width=.4pt,fill=white,rounded corners=1] (748,652) rectangle (756,658);\end{scope}
\begin{scope}[scale=.6666667,draw=black,text=black]\path[draw=black,line width=.4pt,fill=white,rounded corners=1] (748,663) rectangle (756,669);\end{scope}
\begin{scope}[scale=.6666667,draw=black,text=black]\path[draw=black,line width=.4pt,fill=white,rounded corners=1] (759,630) rectangle (767,636);\end{scope}
\begin{scope}[scale=.6666667,draw=black,text=black]\path[draw=black,line width=.4pt,fill=white,rounded corners=1] (759,641) rectangle (767,647);\end{scope}
\begin{scope}[scale=.6666667,draw=black,text=black]\path[draw=black,line width=.4pt,fill=white,rounded corners=1] (759,652) rectangle (767,658);\end{scope}
\begin{scope}[scale=.6666667,draw=black,text=black]\path[draw=black,line width=.4pt,fill=white,rounded corners=1] (759,663) rectangle (767,669);\end{scope}
\q{153}{708}{155}{13}{2元5角}
\q{371}{708}{155}{13}{7角}
\q{533}{708}{155}{13}{6元8分}
\q{704}{708}{155}{13}{35元5角}
\q{143}{756}{28}{13}{（}
\q{214}{756}{25}{13}{）}
\q{234}{756}{45}{13}{元}
\q{337}{756}{28}{13}{（}
\q{405}{756}{25}{13}{）}
\q{425}{756}{45}{13}{元}
\q{524}{756}{28}{13}{（}
\q{590}{756}{25}{13}{）}
\q{610}{756}{45}{13}{元}
\q{699}{756}{28}{13}{（}
\q{768}{756}{25}{13}{）}
\q{788}{756}{45}{13}{元}
\hdsection{4}{816}{解决问题。}
\q{109}{866}{734}{12.5}{6名同学参加50米短跑比赛，成绩如下：}
\begin{scope}[scale=.6666667,draw=black,text=black]\path[pattern=north east lines,pattern color=black] (190,909) rectangle (317,987);\end{scope}
\begin{scope}[scale=.6666667,draw=black,text=black]\draw[draw=black,line width=.5pt] (190.000,909.000)--(190.000,987.000);\end{scope}
\begin{scope}[scale=.6666667,draw=black,text=black]\draw[draw=black,line width=.5pt] (317.000,909.000)--(317.000,987.000);\end{scope}
\begin{scope}[scale=.6666667,draw=black,text=black]\draw[draw=black,line width=.5pt] (395.000,909.000)--(395.000,987.000);\end{scope}
\begin{scope}[scale=.6666667,draw=black,text=black]\draw[draw=black,line width=.5pt] (472.000,909.000)--(472.000,987.000);\end{scope}
\begin{scope}[scale=.6666667,draw=black,text=black]\draw[draw=black,line width=.5pt] (547.000,909.000)--(547.000,987.000);\end{scope}
\begin{scope}[scale=.6666667,draw=black,text=black]\draw[draw=black,line width=.5pt] (623.000,909.000)--(623.000,987.000);\end{scope}
\begin{scope}[scale=.6666667,draw=black,text=black]\draw[draw=black,line width=.5pt] (699.000,909.000)--(699.000,987.000);\end{scope}
\begin{scope}[scale=.6666667,draw=black,text=black]\draw[draw=black,line width=.5pt] (774.000,909.000)--(774.000,987.000);\end{scope}
\begin{scope}[scale=.6666667,draw=black,text=black]\draw[draw=black,line width=.5pt] (190.000,909.000)--(774.000,909.000);\end{scope}
\begin{scope}[scale=.6666667,draw=black,text=black]\draw[draw=black,line width=.5pt] (190.000,948.000)--(774.000,948.000);\end{scope}
\begin{scope}[scale=.6666667,draw=black,text=black]\draw[draw=black,line width=.5pt] (190.000,987.000)--(774.000,987.000);\end{scope}
\begin{scope}[scale=.6666667,draw=black,text=black]\node[inner sep=0pt,font=\normalsize,text=black] at (253.500,928.500) {姓名};\end{scope}
\begin{scope}[scale=.6666667,draw=black,text=black]\node[inner sep=0pt,font=\normalsize,text=black] at (356.000,928.500) {小强};\end{scope}
\begin{scope}[scale=.6666667,draw=black,text=black]\node[inner sep=0pt,font=\normalsize,text=black] at (433.500,928.500) {小东};\end{scope}
\begin{scope}[scale=.6666667,draw=black,text=black]\node[inner sep=0pt,font=\normalsize,text=black] at (509.500,928.500) {小刚};\end{scope}
\begin{scope}[scale=.6666667,draw=black,text=black]\node[inner sep=0pt,font=\normalsize,text=black] at (585.000,928.500) {小庆};\end{scope}
\begin{scope}[scale=.6666667,draw=black,text=black]\node[inner sep=0pt,font=\normalsize,text=black] at (661.000,928.500) {小杰};\end{scope}
\begin{scope}[scale=.6666667,draw=black,text=black]\node[inner sep=0pt,font=\normalsize,text=black] at (736.500,928.500) {小涛};\end{scope}
\begin{scope}[scale=.6666667,draw=black,text=black]\node[inner sep=0pt,font=\normalsize,text=black] at (253.500,967.500) {成绩/秒};\end{scope}
\begin{scope}[scale=.6666667,draw=black,text=black]\node[inner sep=0pt,font=\normalsize,text=black] at (356.000,967.500) {8.9};\end{scope}
\begin{scope}[scale=.6666667,draw=black,text=black]\node[inner sep=0pt,font=\normalsize,text=black] at (433.500,967.500) {8.1};\end{scope}
\begin{scope}[scale=.6666667,draw=black,text=black]\node[inner sep=0pt,font=\normalsize,text=black] at (509.500,967.500) {9.6};\end{scope}
\begin{scope}[scale=.6666667,draw=black,text=black]\node[inner sep=0pt,font=\normalsize,text=black] at (585.000,967.500) {8.3};\end{scope}
\begin{scope}[scale=.6666667,draw=black,text=black]\node[inner sep=0pt,font=\normalsize,text=black] at (661.000,967.500) {9.0};\end{scope}
\begin{scope}[scale=.6666667,draw=black,text=black]\node[inner sep=0pt,font=\normalsize,text=black] at (736.500,967.500) {8.7};\end{scope}
\q{109}{1012}{700}{12.5}{请把前三名的名字写在领奖台上。}
\begin{scope}[scale=.6666667,draw=black,text=black]\draw[draw=black,line width=.6pt,fill=bookpink] (397.000,1140.000)--(420.000,1127.000)--(441.000,1127.000)--(441.000,1140.000)--cycle;\end{scope}
\begin{scope}[scale=.6666667,draw=black,text=black]\draw[draw=black,line width=.6pt,fill=bookpink] (441.000,1116.000)--(465.000,1102.000)--(510.000,1102.000)--(486.000,1116.000)--cycle;\end{scope}
\begin{scope}[scale=.6666667,draw=black,text=black]\draw[draw=black,line width=.6pt,fill=bookpink] (486.000,1148.000)--(510.000,1134.000)--(554.000,1134.000)--(529.000,1148.000)--cycle;\end{scope}
\begin{scope}[scale=.6666667,draw=black,text=black]\draw[draw=black,line width=.6pt,fill=bookpink!75!black] (486.000,1116.000)--(510.000,1102.000)--(510.000,1134.000)--(486.000,1148.000)--cycle;\end{scope}
\begin{scope}[scale=.6666667,draw=black,text=black]\draw[draw=black,line width=.6pt,fill=bookpink!75!black] (529.000,1148.000)--(554.000,1134.000)--(554.000,1157.000)--(529.000,1171.000)--cycle;\end{scope}
\begin{scope}[scale=.6666667,draw=black,text=black]\draw[draw=black,line width=.6pt,fill=bookpink] (397.000,1140.000)--(441.000,1140.000)--(441.000,1116.000)--(486.000,1116.000)--(486.000,1148.000)--(529.000,1148.000)--(529.000,1171.000)--(397.000,1171.000)--cycle;\end{scope}
\begin{scope}[scale=.6666667,draw=black,text=black]\node[inner sep=0pt,font=\normalsize,text=black] at (464.000,1133.000) {1};\end{scope}
\begin{scope}[scale=.6666667,draw=black,text=black]\node[inner sep=0pt,font=\normalsize,text=black] at (417.000,1156.000) {2};\end{scope}
\begin{scope}[scale=.6666667,draw=black,text=black]\node[inner sep=0pt,font=\normalsize,text=black] at (506.000,1159.000) {3};\end{scope}

\begin{scope}[color=answerred,draw=answerred,text=answerred]\ans{114.00}{134.00}{63.33}{9.5}{$\vphantom{0}1.6$}
\ans{121.33}{164.00}{63.33}{9.5}{$\vphantom{0}5.05$}
\ans{130.67}{194.00}{63.33}{9.5}{$\vphantom{0}70.04$}
\ans{247.33}{134.67}{63.33}{9.5}{$\vphantom{0}0.8$}
\ans{254.67}{164.00}{63.33}{9.5}{$\vphantom{0}0.9$}
\ans{254.67}{194.00}{63.33}{9.5}{$\vphantom{0}0.09$}
\ans{391.33}{135.33}{63.33}{9.5}{$\vphantom{0}10.06$}
\ans{382.00}{164.67}{63.33}{9.5}{$\vphantom{0}30.7$}
\ans{391.33}{194.67}{63.33}{9.5}{$\vphantom{0}13.09$}
\ans{512.67}{134.67}{63.33}{9.5}{$\vphantom{0}80$}
\ans{523.33}{164.67}{63.33}{9.5}{$\vphantom{0}0.032$}
\ans{523.33}{194.67}{63.33}{9.5}{$\vphantom{0}0.704$}
\ans{118.00}{264.67}{66.67}{10}{$\vphantom{0}10.00$}
\ans{309.33}{264.67}{66.67}{10}{$\vphantom{0}3.20$}
\ans{480.67}{264.67}{66.67}{10}{$\vphantom{0}21.80$}
\ans{122.00}{295.33}{66.67}{10}{$\vphantom{0}0.70$}
\ans{294.67}{295.33}{66.67}{10}{$\vphantom{0}1.70$}
\ans{494.67}{295.33}{66.67}{10}{$\vphantom{0}20.70$}
\ans{129.33}{326.00}{66.67}{10}{$\vphantom{0}10.60$}
\ans{294.67}{326.00}{66.67}{10}{$\vphantom{0}6.90$}
\ans{488.67}{326.00}{66.67}{10}{$\vphantom{0}1.05$}
\ans{109.33}{503.33}{58.67}{11}{$\vphantom{0}2.50$}
\ans{238.00}{503.33}{58.67}{11}{$\vphantom{0}0.70$}
\ans{362.67}{503.33}{58.67}{11}{$\vphantom{0}6.08$}
\ans{479.33}{503.33}{58.67}{11}{$\vphantom{0}35.50$}
\ans{72.67}{696.67}{467.33}{10}{用时越少，名次越靠前：$8.1<8.3<8.7<8.9<9.0<9.6$。}
\ans{294.00}{719.33}{58.67}{10}{小东}
\ans{244.00}{748.00}{48.00}{9.5}{小庆}
\ans{363.33}{753.33}{48.00}{9.5}{小涛}
\draw[answerred,line width=.7pt] (270.000,757.333) -- (280.000,763.333);
\draw[answerred,line width=.7pt] (365.333,764.667) -- (346.667,766.667);
\ans{110.00}{525.00}{445}{9.5}{第3题单位换算的完整过程见本页后的续页。}
\ans{72.67}{792}{470}{10}{答：第一名是小东，第二名是小庆，第三名是小涛。}
\end{scope}
\end{scope}\pgfresetboundingbox\path[use as bounding box] (230.0,715.0) rectangle (420.0,789.0);
\end{tikzpicture}}\par
\Answer{答：第一名是小东，第二名是小庆，第三名是小涛。}



\QHeading{第4题　短跑名次的另一种确定方法}

\Answer{不用一次排完6名，也可以先选最短时间8.1秒为第一名，再在剩下的人中选8.3秒为第二名，最后选8.7秒为第三名。对应姓名仍是小东、小庆、小涛。短跑按用时从少到多排名，不是分数越高越好。}
